\chapter{From residual to score: residual Mahalanobis, the ring and integral images}\label{ch:residualmaha}

\begin{onething}
L1, SAM and \texttt{combined} measure how \emph{big} a pixel's reconstruction error is, but in a real scene
most big errors are shared by thousands of ordinary pixels. Residual Mahalanobis learns the directions those
ordinary errors take, divides each by its ordinary spread, and measures what is left from the average of a
ring of neighbours; an error then scores high only if it is \emph{unusual}, however small it is, and
integral images make the ring cost four look-ups per box instead of thousands of additions.
\end{onething}

\Cref{ch:inference} ended with a reconstruction $\xhat$ of the whole scene, every pixel predicted while it
was hidden. Subtracting gives the \emph{residual}\index{residual} at every pixel,
\begin{equation}
\vect{r}=\vect{x}-\xhat\in\R^{C},
\end{equation}
one number per band: $C=165$ on the common grid, about 185 for a native-band model. A decision needs one
number per pixel. How you collapse $C$ numbers into one decides what you find. This chapter builds the
collapse that has been Indradhanu's default since 2026-10-15, \code{residual_mahalanobis}, in
\file{app/utils/anomaly_detection/residual_mahalanobis.py} (112 lines). It uses the Mahalanobis distance of
\cref{ch:rx}, the covariance and eigenvectors of \cref{ch:linalg}, and the local windows of
\cref{ch:localrx}, applied to a new input: the model's mistakes rather than the scene itself.

%======================================================================
\section{Size scores, and the big errors that are not anomalies}\label{sec:residualmaha:size}

\subsection{Four ways to measure a residual's size}

Until 2026-10-15 every score the product offered for Indradhanu measured the \emph{size} of $\vect{r}$
(\file{app/utils/anomaly_detection/scoring.py}, lines 95--107):
\begin{align*}
\text{L1}(\vect{r}) &= \frac1C\sum_{c=1}^{C}\abs{r_c} &&\text{mean absolute error},\\
\text{MSE}(\vect{r}) &= \frac1C\sum_{c=1}^{C}r_c^2 &&\text{mean squared error},\\
\SAM(\vect{x},\xhat) &= \arccos\frac{\ip{\vect{x}}{\xhat}}{\norm{\vect{x}}\,\norm{\xhat}} &&\text{angle between spectra},\\
\text{combined} &= 0.5\,\frac{\text{L1}}{\max\text{L1}}+0.5\,\frac{\SAM}{\max\SAM} &&\text{both, maxima over valid pixels}.
\end{align*}
\Cref{ch:losses} worked L1 and SAM by hand and showed that they find different things: L1 flags a pixel that
is uniformly too bright, SAM one whose spectral \emph{shape} is wrong. Both are sizes. Size is the right
question only if every large error is rare.

\subsection{Three kinds of large error that are not rare}

On the four labelled PRISMA scenes the score maps showed three kinds of large error, each shared by thousands
of ordinary pixels (IW45):
\begin{itemize}
\item \textbf{Detector stripes.} A push-broom spectrometer such as PRISMA records each image column with its
  own detector element. One element calibrated slightly off makes every pixel in its column err the same
  way, and the score map shows a stripe. On set~4 the \texttt{combined} map was dominated by them.
\item \textbf{Haze.} A thin layer of atmosphere over part of a scene lifts the blue end of every spectrum
  beneath it. The model, trained mostly on clear ground, under-predicts all those pixels by a similar
  amount. Set~4 again.
\item \textbf{Dark sea.} On set~1, three quarters sea, the map was dominated by the sea, through the trap
  worked next.
\end{itemize}
In each case thousands of pixels share one \emph{kind} of error. A size score ranks all of them above a small
real target. The model is not at fault; the score asks the wrong question.

\subsection{The dark-sea trap, and why its patch is a patch}

SAM is an angle, so it measures error \emph{relative to} the signal. Give a dark sea pixel and a bright land
pixel exactly the same residual, $(+0.01,-0.01,+0.01)$ (AD12):
\begin{center}\small
\begin{tabular}{@{}lccccc@{}}
\toprule
 & band 1 & band 2 & band 3 & L1 & SAM\\
\midrule
sea, true & 0.02 & 0.03 & 0.01 & & \\
sea, reconstruction & 0.03 & 0.02 & 0.02 & 0.010 & 0.4335 rad\\
land, true & 0.20 & 0.30 & 0.10 & & \\
land, reconstruction & 0.21 & 0.29 & 0.11 & 0.010 & 0.0463 rad\\
\bottomrule
\end{tabular}
\end{center}
L1 cannot tell them apart. SAM gives the sea pixel nine times the angle. Measured on set~1 with the F001
model, the median angle is 0.34~rad on sea, 0.13 on land and 0.18 on the labelled
anomalies \tagM{docs/05-detectors.md}: ordinary sea outranks the targets.

\code{compute_score(..., sam_dark_fade=t)} multiplies the angle by $w=\operatorname{clip}(\bar{x}/t,0,1)$,
where $\bar{x}$ is the pixel's mean reflectance (\file{scoring.py}, line 131). With $t=0.05$:
\begin{align*}
\text{sea, } \bar x = 0.02:\quad w &= \operatorname{clip}(0.02/0.05,0,1) = \operatorname{clip}(0.4,0,1) = 0.4,\\
\text{land, } \bar x = 0.20:\quad w &= \operatorname{clip}(0.20/0.05,0,1) = \operatorname{clip}(4.0,0,1) = 1.0.
\end{align*}
\begin{sloppypar}
It works on set~1 (F001 at epoch 52: $0.731\to0.887$ \tagM{indradhanu workings/46-the-ring-and-the-comparison.md})
and leaves the other scenes almost untouched. That is the problem. It is a patch: only set~1 has water,
so four scenes cannot tell whether $0.05$ is right, and the result is steep in $t$ (F001 at epoch 14 scores
$0.524$ on set~1 with $t=0.03$ and $0.810$ with $t=0.08$ \tagM{indradhanu workings/45-scoring-the-direction.md}).
It fixes darkness only; stripes and haze are untouched. And the next failure would need another special case.
The fade is opt-in (\code{None} by default) and the product does not pass it.
\end{sloppypar}

%======================================================================
\section{A residual is a direction}\label{sec:residualmaha:direction}

Think of a residual as an arrow in a $C$-dimensional space, one axis per band. Its \emph{length} is roughly
what L1 measures. Its \emph{direction} says which bands erred, and how they erred together.

Now look at what the shared errors have in common. Haze pushes many bands the same way at once: one
direction. A miscalibrated detector column pushes its pixels another way: one direction. Darkness of the sea:
another. A scene's ordinary errors are large along \emph{a few shared directions} and tiny along all the
others. A real anomaly, a material the model never learned to predict there, errs in a direction the
background rarely takes, even if the arrow is short.

\begin{analogy}
A doctor reads twenty blood tests. A dehydrated patient has every concentration a little high, but they all
rise \emph{together}, because there is less water; the doctor reads them as one ordinary finding, not twenty
alarms. One test far outside \emph{its own} normal range while the rest are ordinary is the result that
matters, even if its deviation is smaller in absolute units than the dehydration's. Residual Mahalanobis
reads residuals the same way: movements many pixels share are one ordinary finding, and the scale of each
test's ``normal range'' is learned from the scene itself. A movement nobody else makes is the alarm.
\end{analogy}

\begin{definition}[Residual covariance and residual Mahalanobis distance]
Let $\vect{r}_1,\dots,\vect{r}_N$ be the residuals of the scene's valid pixels. Their mean is
$\muv_r=\frac1N\sum_i\vect{r}_i$ and their \emph{residual covariance}\index{residual covariance} is the
$C\times C$ matrix
\begin{equation}
\Cov_r=\frac{1}{N-1}\sum_{i=1}^{N}(\vect{r}_i-\muv_r)(\vect{r}_i-\muv_r)\T.
\label{eq:residualmaha:cov}
\end{equation}
The \emph{residual Mahalanobis distance}\index{residual Mahalanobis distance} of a pixel with residual
$\vect{r}$ is
\begin{equation}
D^2(\vect{r})=(\vect{r}-\muv_r)\T\,\Cov_r\inv\,(\vect{r}-\muv_r).
\label{eq:residualmaha:d2}
\end{equation}
\end{definition}
The divisor $N-1$ is the one \code{np.cov} uses, which the code calls; \cref{ch:linalg} explains why it is
the unbiased choice. \Cref{eq:residualmaha:d2} is exactly the RX detector of \cref{ch:rx} with the cube
replaced by the residual. The rest of this section computes it by hand.

\subsection{Worked by hand: three bands, four ordinary pixels}

\begin{worked}[label={we:residualmaha:three}]{Residual Mahalanobis on three bands}
Do this on paper; it takes about fifteen minutes. Three bands A, B and C. Four ordinary pixels have these
residuals (in units of $0.001$ reflectance, to keep the numbers whole):
\begin{center}
\begin{tabular}{@{}cccc@{}}
\toprule
pixel & $r_A$ & $r_B$ & $r_C$\\
\midrule
1 & 1 & 2 & $-1$\\
2 & 1 & $-1$ & $-1$\\
3 & $-2$ & $-1$ & $-1$\\
4 & 0 & 0 & 3\\
\bottomrule
\end{tabular}
\end{center}
Three test pixels: $\vect{P}=(3,3,0)$, $\vect{Q}=(2,-2,0)$ and $\vect{R}=(0,0,4)$.

\textbf{Step 1: the mean, one band at a time.}
\begin{align*}
\mu_A &= \frac{1+1+(-2)+0}{4} = \frac{0}{4} = 0,\\
\mu_B &= \frac{2+(-1)+(-1)+0}{4} = \frac{0}{4} = 0,\\
\mu_C &= \frac{(-1)+(-1)+(-1)+3}{4} = \frac{0}{4} = 0.
\end{align*}
The mean residual is $(0,0,0)$, so every deviation from the mean is the residual itself.

\textbf{Step 2: the covariance, one entry at a time.} Each entry is a sum of products over the four pixels,
divided by $N-1=4-1=3$. The three variances:
\begin{align*}
\sigma_{AA} &= \frac{1^2+1^2+(-2)^2+0^2}{3} = \frac{1+1+4+0}{3} = \frac{6}{3} = 2,\\
\sigma_{BB} &= \frac{2^2+(-1)^2+(-1)^2+0^2}{3} = \frac{4+1+1+0}{3} = \frac{6}{3} = 2,\\
\sigma_{CC} &= \frac{(-1)^2+(-1)^2+(-1)^2+3^2}{3} = \frac{1+1+1+9}{3} = \frac{12}{3} = 4.
\end{align*}
The three covariances:
\begin{align*}
\sigma_{AB} &= \frac{1\cdot2+1\cdot(-1)+(-2)(-1)+0\cdot0}{3} = \frac{2-1+2+0}{3} = \frac{3}{3} = 1,\\
\sigma_{AC} &= \frac{1\cdot(-1)+1\cdot(-1)+(-2)(-1)+0\cdot3}{3} = \frac{-1-1+2+0}{3} = \frac{0}{3} = 0,\\
\sigma_{BC} &= \frac{2\cdot(-1)+(-1)(-1)+(-1)(-1)+0\cdot3}{3} = \frac{-2+1+1+0}{3} = \frac{0}{3} = 0.
\end{align*}
A covariance matrix is symmetric ($\sigma_{BA}=\sigma_{AB}$), so
\[
\Cov_r=\begin{pmatrix}2&1&0\\1&2&0\\0&0&4\end{pmatrix}.
\]
Read it: bands A and B each vary by 2 and move together ($\sigma_{AB}=1>0$); band C varies twice as much and
independently of the other two.

\textbf{Step 3: the inverse.} The matrix is \emph{block diagonal}: an A--B block and a C block, with zeros
between them. Its inverse is the inverse of each block, in place. The C block is the number 4, whose inverse
is $\tfrac14$. For the A--B block use the rule for a $2\times2$ matrix,
\begin{equation}
\begin{pmatrix}a&b\\c&d\end{pmatrix}\inv=\frac{1}{ad-bc}\begin{pmatrix}d&-b\\-c&a\end{pmatrix},
\label{eq:residualmaha:inv2}
\end{equation}
which you can check by multiplying: the top-left entry of the product is
$\frac{1}{ad-bc}(a\cdot d+b\cdot(-c))=\frac{ad-bc}{ad-bc}=1$, and the top-right is
$\frac{1}{ad-bc}(a\cdot(-b)+b\cdot a)=0$. Here $a=2$, $b=1$, $c=1$, $d=2$:
\begin{align*}
ad-bc &= 2\cdot2-1\cdot1 = 4-1 = 3,\\
\begin{pmatrix}2&1\\1&2\end{pmatrix}\inv &= \frac13\begin{pmatrix}2&-1\\-1&2\end{pmatrix}.
\end{align*}
So
\[
\Cov_r\inv=\begin{pmatrix}\tfrac23&-\tfrac13&0\\[2pt]-\tfrac13&\tfrac23&0\\[2pt]0&0&\tfrac14\end{pmatrix}.
\]
Check one row of $\Cov_r\Cov_r\inv$ against the identity. Row 1 of $\Cov_r$ is $(2,1,0)$:
\begin{align*}
\text{with column 1: } & 2\cdot\tfrac23+1\cdot(-\tfrac13)+0\cdot0 = \tfrac43-\tfrac13 = 1,\\
\text{with column 2: } & 2\cdot(-\tfrac13)+1\cdot\tfrac23+0\cdot0 = -\tfrac23+\tfrac23 = 0,\\
\text{with column 3: } & 2\cdot0+1\cdot0+0\cdot\tfrac14 = 0.
\end{align*}
That is $(1,0,0)$, the first row of the identity.

\textbf{Step 4: score each test pixel.} First multiply $\Cov_r\inv$ by the residual, then take the dot product
with the residual. For $\vect{P}=(3,3,0)$:
\begin{align*}
\Cov_r\inv\vect{P} &= \left(\tfrac23\cdot3-\tfrac13\cdot3+0,\ -\tfrac13\cdot3+\tfrac23\cdot3+0,\ \tfrac14\cdot0\right)\\
&= (2-1,\ -1+2,\ 0) = (1,1,0),\\
D^2(\vect{P}) &= 3\cdot1+3\cdot1+0\cdot0 = 6.
\end{align*}
For $\vect{Q}=(2,-2,0)$:
\begin{align*}
\Cov_r\inv\vect{Q} &= \left(\tfrac23\cdot2-\tfrac13\cdot(-2),\ -\tfrac13\cdot2+\tfrac23\cdot(-2),\ 0\right)\\
&= \left(\tfrac43+\tfrac23,\ -\tfrac23-\tfrac43,\ 0\right) = (2,-2,0),\\
D^2(\vect{Q}) &= 2\cdot2+(-2)(-2)+0 = 4+4 = 8.
\end{align*}
For $\vect{R}=(0,0,4)$:
\begin{align*}
\Cov_r\inv\vect{R} &= (0,\ 0,\ \tfrac14\cdot4) = (0,0,1),\\
D^2(\vect{R}) &= 0+0+4\cdot1 = 4.
\end{align*}
And one ordinary pixel, pixel~1, $(1,2,-1)$:
\begin{align*}
\Cov_r\inv\vect{r}_1 &= \left(\tfrac23\cdot1-\tfrac13\cdot2,\ -\tfrac13\cdot1+\tfrac23\cdot2,\ \tfrac14\cdot(-1)\right)
 = \left(0,\ 1,\ -\tfrac14\right),\\
D^2(\vect{r}_1) &= 1\cdot0+2\cdot1+(-1)(-\tfrac14) = 0+2+\tfrac14 = 2.25.
\end{align*}
\end{worked}

Put the answers side by side (all computed above or in the book's script \tagC):
\begin{center}\small
\begin{tabular}{@{}lcccc@{}}
\toprule
pixel & residual & length $\norm{\vect{r}}$ & L1 & $D^2$\\
\midrule
ordinary 1--4 & & 1.73--3.00 & 1.00--1.33 & 2.25 each\\
$\vect{P}$ & $(3,3,0)$ & 4.24 & 2.00 & 6\\
$\vect{R}$ & $(0,0,4)$ & 4.00 & 1.33 & 4\\
$\vect{Q}$ & $(2,-2,0)$ & 2.83 & 1.33 & \textbf{8}\\
\bottomrule
\end{tabular}
\end{center}
\textbf{The ranking flips.} By length, $\vect{P}>\vect{R}>\vect{Q}$; by L1, $\vect{P}$ first. By $D^2$,
$\vect{Q}$ is the most anomalous and $\vect{P}$ only second. $\vect{P}$ errs along ``A and B together'', the
direction the ordinary pixels share. $\vect{R}$ errs in band C, which is noisy anyway. $\vect{Q}$ is the
shortest arrow, but it has A and B \emph{disagreeing}, which the ordinary pixels rarely do.

A curiosity you may have noticed: all four ordinary pixels score exactly $2.25$. That is not luck. With
$N$ pixels in $N-1$ dimensions, every pixel used to estimate the covariance scores
$(N-1)^2/N=3^2/4=2.25$ against it. The covariance has used up all its freedom fitting those few pixels. The
code avoids this regime by estimating the covariance from up to 200,000 pixels.

%======================================================================
\section{Whitening: rotate, divide, measure}\label{sec:residualmaha:whiten}

\subsection{The eigen-decomposition turns the inverse into a sum}

The code never forms $\Cov_r\inv$. It uses the eigenvectors (\cref{ch:linalg}). A symmetric matrix has an
\emph{eigen-decomposition}\index{eigen-decomposition}
\begin{equation}
\Cov_r=\mat{U}\boldsymbol{\Lambda}\mat{U}\T,
\end{equation}
where the columns $\vect{u}_k$ of $\mat{U}$ are unit-length, mutually perpendicular eigenvectors
($\mat{U}\T\mat{U}=\mat{I}$) and $\boldsymbol{\Lambda}=\diag(\lambda_1,\dots,\lambda_C)$ holds the
eigenvalues, the variance of the residuals along each $\vect{u}_k$. Derive the inverse:
\begin{align*}
\Cov_r\inv &= (\mat{U}\boldsymbol{\Lambda}\mat{U}\T)\inv
\intertext{The inverse of a product is the product of the inverses in reverse order:}
&= (\mat{U}\T)\inv\,\boldsymbol{\Lambda}\inv\,\mat{U}\inv
\intertext{and since $\mat{U}\T\mat{U}=\mat{I}$, the inverse of $\mat{U}$ is $\mat{U}\T$ and the inverse of
$\mat{U}\T$ is $\mat{U}$:}
&= \mat{U}\,\boldsymbol{\Lambda}\inv\,\mat{U}\T.
\end{align*}
Substitute into \cref{eq:residualmaha:d2}, writing $\vect{d}=\vect{r}-\muv_r$:
\begin{align*}
D^2 &= \vect{d}\T\mat{U}\boldsymbol{\Lambda}\inv\mat{U}\T\vect{d}
\intertext{Group the product as $(\mat{U}\T\vect{d})\T\,\boldsymbol{\Lambda}\inv\,(\mat{U}\T\vect{d})$, using
$(\mat{U}\T\vect{d})\T=\vect{d}\T\mat{U}$ (the transpose of a product reverses the order). The vector
$\vect{y}=\mat{U}\T\vect{d}$ has entries $y_k=\vect{u}_k\T\vect{d}$, the coordinates of $\vect{d}$ along each
eigenvector. $\boldsymbol{\Lambda}\inv$ is diagonal, so the product is a sum:}
&= \sum_{k=1}^{C}\frac{y_k^2}{\lambda_k}
 = \sum_{k=1}^{C}\left(\frac{\vect{u}_k\T(\vect{r}-\muv_r)}{\sqrt{\lambda_k}}\right)^2.
\end{align*}
Define the \emph{whitened coordinates}\index{whitening}
\begin{equation}
z_k=\frac{\vect{u}_k\T(\vect{r}-\muv_r)}{\sqrt{\lambda_k}},\qquad D^2=\sum_k z_k^2=\norm{\vect{z}}^2.
\label{eq:residualmaha:whiten}
\end{equation}
Three operations: subtract the mean, \emph{rotate} onto the eigenvectors, \emph{divide} each coordinate by its
direction's standard deviation $\sqrt{\lambda_k}$. After that the ordinary residuals spread equally in every
direction, and the plain squared length is the Mahalanobis distance. Dividing by $\sqrt{\lambda_k}$ and not
by $\lambda_k$ matters: $\lambda_k$ is a variance, in squared units, and a coordinate must be divided by a
spread in its own units.

\begin{worked}[label={we:residualmaha:eigen}]{The same scores by rotate and divide}
The eigenvectors of $\Cov_r$ from \cref{we:residualmaha:three} can be found by inspection. Check each by
multiplying:
\begin{align*}
\vect{u}_1 &= \tfrac{1}{\sqrt2}(1,1,0):& \Cov_r\vect{u}_1 &= \tfrac{1}{\sqrt2}(2+1,\ 1+2,\ 0) = 3\,\vect{u}_1, & \lambda_1 &= 3,\\
\vect{u}_2 &= \tfrac{1}{\sqrt2}(1,-1,0):& \Cov_r\vect{u}_2 &= \tfrac{1}{\sqrt2}(2-1,\ 1-2,\ 0) = 1\,\vect{u}_2, & \lambda_2 &= 1,\\
\vect{u}_3 &= (0,0,1):& \Cov_r\vect{u}_3 &= (0,\ 0,\ 4) = 4\,\vect{u}_3, & \lambda_3 &= 4.
\end{align*}
``A and B together'' spreads with variance 3; ``A and B disagreeing'' with variance 1; band C with variance 4.
Now rotate and divide, with $\tfrac{1}{\sqrt2}\approx0.7071$. For $\vect{P}=(3,3,0)$:
\begin{align*}
y_1 &= 0.7071\,(3+3) = 4.243, & z_1 &= 4.243/\sqrt3 = 4.243/1.732 = 2.449,\\
y_2 &= 0.7071\,(3-3) = 0, & z_2 &= 0/1 = 0,\\
y_3 &= 0, & z_3 &= 0/2 = 0,\\
D^2 &= 2.449^2+0+0 = 6.00.
\end{align*}
For $\vect{Q}=(2,-2,0)$:
\begin{align*}
y_1 &= 0.7071\,(2-2) = 0, & z_1 &= 0,\\
y_2 &= 0.7071\,(2+2) = 2.828, & z_2 &= 2.828/1 = 2.828,\\
y_3 &= 0, & z_3 &= 0,\\
D^2 &= 0+2.828^2+0 = 8.00.
\end{align*}
For $\vect{R}=(0,0,4)$: $y_3=4$, $z_3=4/\sqrt4=4/2=2$, $D^2=2^2=4$. The same three numbers as with the
inverse matrix, as the derivation promised. Running the product's own \code{whitened_pca} on the four ordinary
pixels gives whitened coordinates whose squares sum to $2.25$ for every pixel; it may return an eigenvector
with the opposite sign, which changes nothing after squaring \tagC.
\end{worked}

\coursefig[0.95\linewidth]{iw-45-whitening}{Rotate, divide, measure (course figure 45.1, the course's two-band
version of the example above). (a) Ordinary residuals (grey) spread along the shared diagonal far more than
across it; one test pixel lies further out along the diagonal, the other across it. (b) The same points turned
onto the eigenvectors. (c) After dividing by each direction's spread the ordinary pixels sit on one circle and
the across-the-diagonal pixel is far outside it.}{fig:residualmaha:whitening}

\coursefig[0.9\linewidth]{ad-12-residual-mahalanobis}{Size versus direction (anomaly course figure 12.3). Left:
ordinary residuals spread widely along direction 1 and barely along direction 2. Right: after whitening the
cloud is round; the pixel far out along the common direction is inside it, the one a short way along the rare
direction is far outside.}{fig:residualmaha:ad12}

\subsection{The same steps in the code}

\code{whitened_pca} (\file{residual_mahalanobis.py}, lines 38--49) does steps 1 to 4 for a whole scene:
\begin{lstlisting}
X = residual[:, keep].T.astype(np.float64)          # (N, C): one row per valid pixel
mu = X.mean(0)                                       # (C,)   step 1: the mean
rng = np.random.default_rng(seed)                    # seed = 0: deterministic
sub = X[rng.choice(len(X), size=min(len(X), 200_000), replace=False)] - mu   # (<=200000, C)
evals, evecs = np.linalg.eigh(np.cov(sub.T))         # (C,), (C, C): steps 2-3 on a sample
order = np.argsort(evals)[::-1][:d]                  # (d,)   the d largest spreads, d = 30
P = ((X - mu) @ evecs[:, order]) / np.sqrt(np.maximum(evals[order], 1e-12))  # (N, d): rotate, divide
out = np.zeros((len(order),) + keep.shape, dtype=np.float32)                 # (d, H, W)
out[:, keep] = P.T                                   # whitened coordinates back on the image grid
\end{lstlisting}
and, with no local window, the score is \code{(X ** 2).sum(0)}: \cref{eq:residualmaha:whiten}. Three
details:
\begin{itemize}
\item \textbf{A sample, with a fixed seed.} The covariance comes from at most 200,000 valid pixels
  (\code{_SAMPLE}, line 35) drawn with seed 0, so the score is deterministic and the cost of the eigen
  problem does not grow with the scene. The mean uses every valid pixel.
\item \textbf{The mean is the scene's own.} Every scene gets its own directions and spreads. A stripe or a
  haze bank that dominates \emph{this} scene is learned, and discounted, here.
\item \textbf{The top $d=30$ directions only} (\code{PCA_COMPONENTS}, line 32), with a floor of $10^{-12}$
  on each eigenvalue. The next section explains both.
\end{itemize}

%======================================================================
\section{Regularisation: why only the top 30 directions}\label{sec:residualmaha:regularise}

\subsection{Small eigenvalues explode}
\index{regularisation}\index{RX detector!plain, dropped}

\Cref{eq:residualmaha:whiten} divides by $\sqrt{\lambda_k}$. If an estimated eigenvalue is tiny, a tiny
residual along it becomes huge. Watch it happen on a $2\times2$ covariance with unit variances and
correlation $c$ close to 1, scoring the small residual $(0.01,-0.01)$, which lies exactly along the
direction of the small eigenvalue $1-c$ (computed by the book's script \tagC):
\begin{center}\small
\begin{tabular}{@{}cccc@{}}
\toprule
$c$ & eigenvalues & diagonal of $\Cov\inv$ & $D^2$ of $(0.01,-0.01)$\\
\midrule
0.9 & 0.1, 1.9 & 5.26 & 0.002\\
0.99 & 0.01, 1.99 & 50.3 & 0.02\\
0.999 & 0.001, 1.999 & 500 & 0.2\\
\bottomrule
\end{tabular}
\end{center}
Every factor of ten in the smallest eigenvalue is a factor of ten in the score. Push it to $\lambda=10^{-10}$
(a standard deviation of $0.00001$) and a residual of $0.001$ along it contributes
\begin{align*}
\frac{0.001^2}{10^{-10}} &= \frac{0.000001}{0.0000000001} = 10{,}000.
\end{align*}
One meaningless direction now outscores everything. With many bands the smallest eigenvalues of an
\emph{estimated} covariance are mostly estimation noise. This is why plain RX, Mahalanobis on the raw cube over
all bands, was dropped from the product on 2026-05-11, with distances reaching about
$10^{11}$ \tagM{CLAUDE.md}.

\subsection{Three remedies, and the one the code uses}

\begin{itemize}
\item \textbf{Truncation}\index{truncation} (what the code does): keep only the $d$ directions with the largest eigenvalues.
  The smallest kept eigenvalue of 30 out of about 185 carries real variance; the explosive ones are
  discarded.
\item \textbf{A floor}: \code{np.maximum(evals, 1e-12)} only prevents a division by exactly zero; at
  $10^{-12}$ it does not tame anything else.
\item \textbf{Shrinkage}\index{shrinkage} (not used here): add a small $\varepsilon$ to every eigenvalue,
  $\Cov+\varepsilon\mat{I}$ (\cref{ch:localrx}). At $c=0.999$ with $\varepsilon=0.01$ the small eigenvalue
  becomes $0.011$ and the score of $(0.01,-0.01)$ drops from $0.2$ to $0.0002/0.011=0.018$.
\end{itemize}

\paragraph{How many numbers does the covariance estimate?} A $d\times d$ symmetric matrix has
$d(d+1)/2$ distinct entries: the diagonal plus the triangle above it. For $d=30$:
\begin{align*}
\frac{30\times31}{2} = \frac{930}{2} = 465,
\end{align*}
from up to 200,000 pixels: well determined. The full covariance over 185 bands would need
$185\times186/2=17{,}205$ entries, and over 165 bands $165\times166/2=13{,}695$; the smallest directions of
those are the ones that explode.

\paragraph{What truncation costs.} Truncation throws away the directions in which the background barely moves,
and those are exactly where a small anomaly like $\vect{Q}$ stands out. Repeat \cref{we:residualmaha:eigen}
keeping only the top $d=2$ directions, $\lambda_3=4$ and $\lambda_1=3$, and dropping $\lambda_2=1$:
\begin{align*}
\vect{P}: & \ z=(0,\ 2.449) \to D^2 = 6.0,\\
\vect{R}: & \ z=(2,\ 0) \to D^2 = 4.0,\\
\vect{Q}: & \ z=(0,\ 0) \to D^2 = 0.
\end{align*}
$\vect{Q}$ disappears. In a real scene the trade is better than in this toy, because the 155 discarded
directions of a 185-band residual are mostly noise and a real material differs from its background in many
bands at once, so it has a component along the kept directions too. But the trade is real: keeping 30 is a
choice for stability over sensitivity in the quietest directions, made once on four labelled scenes.

\subsection{Why on the residual, and how much the model adds}

RX on the cube, Reed and Yu's 1990 detector \citep{reed1990adaptive}, is the same distance with a different input. A raw scene mixes very different surfaces (sea,
crops, towns), so one set of directions fits it badly. The model has already explained ``what surface is
this'', and what it leaves behind is closer to one cloud of ordinary errors, which is what a single
Mahalanobis distance assumes. Be honest about the size of that effect. On the four PRISMA scenes with the
production model v0.2.0 (mean ROC-AUC, \cref{ch:measures}):
\begin{center}\small
\begin{tabular}{@{}lc@{}}
\toprule
score & mean ROC-AUC\\
\midrule
L1 on the residual & 0.702\\
Mahalanobis, top 30, on the \emph{cube} & 0.831\\
Mahalanobis, top 30, on the \emph{residual} (global mean) & 0.839\\
\bottomrule
\end{tabular}

\smallskip
Source: \tagM{research/experiments/dual-window-mahalanobis/README.md}
\end{center}
\begin{align*}
\text{gain from scoring the direction:} &\quad 0.839-0.702 = 0.137,\\
\text{gain from using the residual:} &\quad 0.839-0.831 = 0.008.
\end{align*}
Most of the improvement comes from whitening, not from the model. The README's own words: the residual
``adds a small, consistent edge on top of a strong classical baseline''.

%======================================================================
\section{The ring: a local mean with a global covariance}\label{sec:residualmaha:ring}

\subsection{The global mean assumes one centre}

\Cref{eq:residualmaha:d2} measures every pixel from the scene's mean residual $\muv_r$. That assumes every
ordinary pixel's error scatters around one centre. It does not quite: the dual-window README puts it as ``the
model has a per-land-cover bias: its residual over water is not centred where its residual over crops is.''
Work it on one whitened coordinate. In a coastal scene the sea pixels sit around $+1.5$ along it and the land
pixels around $0$ \tagA:
\begin{center}\small
\begin{tabular}{@{}lcccc@{}}
\toprule
pixel & value & global score & ring mean & local score\\
\midrule
ordinary sea & 1.8 & $1.8^2=3.24$ & 1.5 & $(1.8-1.5)^2=0.3^2=0.09$\\
land target & 1.6 & $1.6^2=2.56$ & 0.0 & $(1.6-0.0)^2=1.6^2=2.56$\\
\bottomrule
\end{tabular}
\end{center}
Measured from the scene's centre the ordinary sea pixel outranks the target. Measured from its own
neighbourhood it drops to almost nothing, and the target stands out.

\subsection{The local score, derived}

Let $\muv_{\text{ring}}(p)$ be the mean residual over a ring of pixels around $p$. The \emph{local-mean
Mahalanobis distance}\index{local-mean Mahalanobis distance} replaces the global mean with it but keeps the
global covariance:
\begin{equation}
D^2_{\text{local}}(p)=(\vect{r}_p-\muv_{\text{ring}}(p))\T\,\Cov_r\inv\,(\vect{r}_p-\muv_{\text{ring}}(p)).
\label{eq:residualmaha:local}
\end{equation}
The code computes it in whitened space. Whitening is a fixed linear map,
$\vect{z}=\mat{W}(\vect{r}-\muv_r)$ with $\mat{W}=\boldsymbol{\Lambda}^{-1/2}\mat{U}\T$. The mean of a linear map
over the ring is the linear map of the mean, because a sum can be taken inside a matrix product:
\begin{align*}
\bar{\vect{z}}_{\text{ring}} &= \frac1n\sum_{q\in\text{ring}}\mat{W}(\vect{r}_q-\muv_r)
 = \mat{W}\left(\frac1n\sum_{q\in\text{ring}}\vect{r}_q-\muv_r\right)
 = \mat{W}(\muv_{\text{ring}}-\muv_r).
\intertext{Subtract it from the pixel's own whitened vector; the global mean cancels:}
\vect{z}_p-\bar{\vect{z}}_{\text{ring}} &= \mat{W}(\vect{r}_p-\muv_r)-\mat{W}(\muv_{\text{ring}}-\muv_r)
 = \mat{W}(\vect{r}_p-\muv_{\text{ring}}).
\intertext{Its squared length is \cref{eq:residualmaha:local}, restricted to the top $d$ directions:}
\norm{\vect{z}_p-\bar{\vect{z}}_{\text{ring}}}^2 &= (\vect{r}_p-\muv_{\text{ring}})\T\mat{W}\T\mat{W}(\vect{r}_p-\muv_{\text{ring}}).
\end{align*}
So the code can whiten once, globally, and then average the whitened maps over rings. The covariance, the
directions and spreads, still come from the whole scene; only the centre is local.

\subsection{The ring and its guard}

\code{LOCAL_WINDOW = (21, 61)} (\file{residual_mahalanobis.py}, line 33) gives two full widths, both odd so
that they centre on the pixel:
\begin{align*}
\text{outer box} &= 61\times61 = 3{,}721 \text{ pixels},\\
\text{guard box} &= 21\times21 = 441 \text{ pixels},\\
\text{ring} &= 3{,}721-441 = 3{,}280 \text{ pixels}.
\end{align*}
The \emph{guard}\index{guard window}\index{ring} is the inner box left out of the average. Without it, a target several
pixels wide would sit partly in its own ``background'', pull the ring mean towards itself and shrink its own
score. A guard 21 wide reaches $\lfloor21/2\rfloor=10$ pixels from the centre, so anything within 10 pixels
(in rows and in columns) is excluded; the outer box reaches $\lfloor61/2\rfloor=30$.

Two edge rules from the code:
\begin{itemize}
\item Only \emph{valid} pixels count: the ring mean is a sum of valid values divided by the number of valid
  pixels.
\item If fewer than \code{MIN_RING_FACTOR} $\times\, d=2\times30=60$ valid pixels remain in a ring (at a scene
  edge, or beside a nodata hole), that pixel falls back to its global-mean score
  (\file{residual_mahalanobis.py}, lines 34 and 67--69).
\end{itemize}

\coursefig[0.9\linewidth]{iw-46-ring}{The ring, and a box sum from four look-ups (course figure 46.1). (a) The
$61\times61$ window minus the hatched $21\times21$ guard leaves a ring of 3,280 pixels whose mean is subtracted
from the centre pixel. (b) A $4\times4$ grid and its zero-led integral image: the shaded $2\times2$ box sums to
6 from four entries of the table.}{fig:residualmaha:ring}

\subsection{Why not a local covariance too?}

That is the classical dual-window RX of \cref{ch:localrx}: mean \emph{and} covariance from the ring. On the
residual it scored worse, 0.814 against the local mean's 0.854 mean AUC with v0.2.0, in step~4 of the
dual-window experiment log \tagM{research/experiments/dual-window-mahalanobis/EXPERIMENT_LOG.md}.
A ring of a few hundred to a few thousand pixels, of which many are correlated neighbours, cannot estimate 465
covariance numbers reliably, and it soaks up nearby targets. The log's own conclusion: ``The local
covariance is the problem, not locality.''

%======================================================================
\section{Summing a ring fast: integral images}\label{sec:residualmaha:integral}

\subsection{The cost of doing it directly}

The ring mean needs, at every pixel, the sum of 3,280 values in each of 30 whitened channels, plus the count of
valid pixels: 31 channels. Done directly on the $1210\times1219$ PRISMA frame of \cref{ch:inference}:
\begin{align*}
\text{pixels} &= 1210\times1219 = 1{,}474{,}990,\\
\text{additions} &= 3{,}280\times31\times1{,}474{,}990\\
&= 101{,}680\times1{,}474{,}990\\
&\approx 1.50\times10^{11}.
\end{align*}
About 150 billion additions \tagC. An integral image reduces each box sum to four look-ups.

\subsection{Definition and construction}

\begin{definition}[Integral image]
The \emph{integral image}\index{integral image} (or \emph{summed-area table}\index{summed-area table}) of an
$H\times W$ grid $v$ is the $(H+1)\times(W+1)$ table
\begin{equation}
I[i][j]=\sum_{r=0}^{i-1}\sum_{c=0}^{j-1}v[r][c],\qquad 0\le i\le H,\ 0\le j\le W,
\end{equation}
the sum of everything strictly above and to the left of position $(i,j)$. Row $0$ and column $0$ are all zero
(empty sums); the table is \emph{zero-led}.
\end{definition}

You can build it entry by entry with a recurrence. The block above-left of $(i+1,j+1)$ is the value $v[i][j]$
plus the block above it plus the block to its left; those two blocks share the block above-left of $(i,j)$,
which has then been counted twice, so subtract it once:
\begin{equation}
I[i+1][j+1]=v[i][j]+I[i][j+1]+I[i+1][j]-I[i][j].
\label{eq:residualmaha:recur}
\end{equation}
The code does something equivalent and faster: a running sum down each column, then a running sum along each
row (\code{cumsum(1).cumsum(2)} on a \shape{(channels, H, W)} array), then a row and column of zeros in front.

\begin{worked}[label={we:residualmaha:integral}]{Building an integral image on a $4\times4$ grid}
Take the course's grid:
\[
v=\begin{pmatrix}1&2&0&1\\3&1&2&0\\0&2&1&3\\1&0&2&1\end{pmatrix}.
\]
\textbf{Running sum down each column.} Each entry is itself plus everything above it. Column 0: $1$, $1+3=4$,
$4+0=4$, $4+1=5$. Column 1: $2$, $2+1=3$, $3+2=5$, $5+0=5$. Column 2: $0$, $0+2=2$, $2+1=3$, $3+2=5$.
Column 3: $1$, $1+0=1$, $1+3=4$, $4+1=5$:
\[
\begin{pmatrix}1&2&0&1\\4&3&2&1\\4&5&3&4\\5&5&5&5\end{pmatrix}.
\]
\textbf{Running sum along each row} of that table. Row 0: $1$, $1+2=3$, $3+0=3$, $3+1=4$. Row 1: $4$,
$4+3=7$, $7+2=9$, $9+1=10$. Row 2: $4$, $4+5=9$, $9+3=12$, $12+4=16$. Row 3: $5$, $5+5=10$, $10+5=15$,
$15+5=20$.

\textbf{Add the zero row and column in front:}
\[
I=\begin{pmatrix}0&0&0&0&0\\0&1&3&3&4\\0&4&7&9&10\\0&4&9&12&16\\0&5&10&15&20\end{pmatrix}.
\]
Check one entry with the recurrence \eqref{eq:residualmaha:recur}, the one for value row 1, column 1:
\begin{align*}
I[2][2] &= v[1][1]+I[1][2]+I[2][1]-I[1][1] = 1+3+4-1 = 7,
\end{align*}
and directly: the top-left $2\times2$ block is $1+2+3+1=7$. The last entry, 20, is the sum of the whole grid.
\end{worked}

\subsection{The four-term box sum}
\index{box sum}\index{inclusion-exclusion}

A box over rows $r_0,\dots,r_1$ and columns $c_0,\dots,c_1$ has sum
\begin{equation}
S=I[r_1+1][c_1+1]-I[r_0][c_1+1]-I[r_1+1][c_0]+I[r_0][c_0].
\label{eq:residualmaha:box}
\end{equation}
Derive it from the definition. Call the box $B$ and name three other regions, each a block that starts at the
top-left corner of the grid:
\begin{itemize}
\item $\text{Big}$: rows $0..r_1$, columns $0..c_1$; its sum is $I[r_1+1][c_1+1]$.
\item $T$ (top strip): rows $0..r_0-1$, columns $0..c_1$; its sum is $I[r_0][c_1+1]$.
\item $L$ (left strip): rows $0..r_1$, columns $0..c_0-1$; its sum is $I[r_1+1][c_0]$.
\end{itemize}
Big is the box plus the parts of $T$ and $L$:
\begin{align*}
\text{Big} &= B + (T\cup L).
\intertext{$T$ and $L$ overlap in the corner $T\cap L$: rows $0..r_0-1$, columns $0..c_0-1$, with sum
$I[r_0][c_0]$. Adding $T$ and $L$ counts the corner twice, so}
\text{sum}(T\cup L) &= \text{sum}(T)+\text{sum}(L)-\text{sum}(T\cap L).
\intertext{Substitute and solve for the box:}
\text{sum}(B) &= \text{sum}(\text{Big})-\text{sum}(T)-\text{sum}(L)+\text{sum}(T\cap L),
\end{align*}
which is \cref{eq:residualmaha:box}. Four look-ups, three additions or subtractions, whatever the size of the
box: a $61\times61$ box costs the same as a $2\times2$.

\begin{worked}[label={we:residualmaha:box}]{A box sum and a ring mean from the table}
With $I$ from \cref{we:residualmaha:integral}. The box of value rows 1--2 and columns 1--2 ($r_0=1$, $r_1=2$,
$c_0=1$, $c_1=2$):
\begin{align*}
S &= I[3][3]-I[1][3]-I[3][1]+I[1][1]\\
&= 12-3-4+1\\
&= 6.
\end{align*}
Check by adding the four values: $1+2+2+1=6$.

\textbf{A ring is one box minus another.} Shrink the windows to a toy: outer box $3\times3$, guard $1\times1$,
around value row 1, column 1. The outer box is rows 0--2, columns 0--2:
\begin{align*}
S_{\text{outer}} &= I[3][3]-I[0][3]-I[3][0]+I[0][0] = 12-0-0+0 = 12.
\end{align*}
The guard is the single value at row 1, column 1:
\begin{align*}
S_{\text{guard}} &= I[2][2]-I[1][2]-I[2][1]+I[1][1] = 7-3-4+1 = 1.
\end{align*}
Ring:
\begin{align*}
\text{ring sum} &= 12-1 = 11,\\
\text{ring count} &= 3\times3-1\times1 = 9-1 = 8,\\
\text{ring mean} &= \frac{11}{8} = 1.375.
\end{align*}
The pixel's value is 1, so its local score in this one-channel toy is $(1-1.375)^2=(-0.375)^2=0.140625$.
Running the product's \code{local_mean_score} on this grid with \code{window=(1, 3)} returns exactly
$0.140625$ at that pixel \tagC.
\end{worked}

\subsection{Counts, padding and the code}

\code{local_mean_score} (\file{residual_mahalanobis.py}, lines 58--69), annotated:
\begin{lstlisting}
d, H, W = X.shape                              # X: (30, H, W) whitened, 0 outside keep
inner, outer = window                          # (21, 61)
r = outer // 2                                 # 30
F = np.concatenate([keep[None], X * keep], 0).astype(np.float64)   # (31, H, W): channel 0 = 1 if valid
I = np.pad(np.pad(F, ((0, 0), (r, r), (r, r))).cumsum(1).cumsum(2),
           ((0, 0), (1, 0), (1, 0)))           # (31, H+2r+1, W+2r+1): zero-led integral image
A = _box(I, outer // 2, r, H, W) - _box(I, inner // 2, r, H, W)   # (31, H, W): ring count, ring sums
n = A[0]                                       # (H, W) valid pixels in each ring
ok = keep & (n >= MIN_RING_FACTOR * d)         # at least 60
s = ((X - A[1:] / np.maximum(n, 1)) ** 2).sum(0)   # (H, W): |z - ring mean|^2
return (np.where(ok, s, (X ** 2).sum(0)) * keep).astype(np.float32)   # fallback: global-mean score
\end{lstlisting}
and \code{_box} (lines 52--55) is \cref{eq:residualmaha:box} for every pixel at once:
\begin{lstlisting}
a, b = r - k, r + k + 1        # k = half-width of the box: 30 for the outer box, 10 for the guard
return I[:, b:b+h, b:b+w] - I[:, a:a+h, b:b+w] - I[:, b:b+h, a:a+w] + I[:, a:a+h, a:a+w]
\end{lstlisting}
Three things make it exact.
\begin{itemize}
\item \textbf{Counts ride along as channel 0.} \code{keep} is 1 for valid pixels, so its box sum is the valid
  count. Invalid pixels were set to 0 in \code{X * keep}, so they add nothing to the sums either.
\item \textbf{Padding by $r$} surrounds the scene with zeros, so a box hanging over the edge counts fewer
  pixels, and the count channel records how many.
\item \textbf{The index arithmetic.} Pixel $(i,j)$ of the scene sits at $(i+r,j+r)$ in the padded array. A box
  of half-width $k$ around it spans padded rows $i+r-k,\dots,i+r+k$, so by \cref{eq:residualmaha:box} it needs
  $I$ at rows $i+(r-k)$ and $i+(r+k+1)$: the slices starting at $a=r-k$ and $b=r+k+1$.
\end{itemize}
The whole score takes about 3--4~seconds per scene on a CPU, against about 46~seconds for F001's
reconstruction \tagM{docs/lld/residual-mahalanobis.md}.

\subsection{Local covariances from the same trick, and why they are not used}

Integral images give any local statistic that is a sum. A local covariance needs the sums of $\vect{x}$ and of
$\vect{x}\vect{x}\T$. Derive it from the definition with divisor $n$:
\begin{align*}
\Cov_{\text{loc}} &= \frac1n\sum_q(\vect{x}_q-\vect{m})(\vect{x}_q-\vect{m})\T
\intertext{Expand the product of the brackets into four terms:}
&= \frac1n\sum_q\vect{x}_q\vect{x}_q\T-\frac1n\sum_q\vect{x}_q\vect{m}\T-\frac1n\sum_q\vect{m}\vect{x}_q\T+\frac1n\sum_q\vect{m}\vect{m}\T
\intertext{and use $\frac1n\sum_q\vect{x}_q=\vect{m}$ in the middle terms and $\frac1n\sum_q1=1$ in the last:}
&= \frac1n\sum_q\vect{x}_q\vect{x}_q\T-\vect{m}\vect{m}\T-\vect{m}\vect{m}\T+\vect{m}\vect{m}\T\\
&= \frac1n\sum_q\vect{x}_q\vect{x}_q\T-\vect{m}\vect{m}\T.
\end{align*}
Each entry $\sum_q x_{q,a}x_{q,b}$ is a box sum of the product image $x_a x_b$, so it too costs four look-ups.

\begin{worked}[label={we:residualmaha:localcov}]{A local covariance from five box sums}
Two channels $a$ and $b$, one $3\times3$ box of $n=9$ pixels (values from the book's script \tagC):
\[
a=\begin{pmatrix}0&0&0\\3&2&2\\2&2&2\end{pmatrix},\qquad b=\begin{pmatrix}0&0&3\\0&1&1\\0&0&0\end{pmatrix}.
\]
The five box sums:
\begin{align*}
\textstyle\sum a &= 0+0+0+3+2+2+2+2+2 = 13,\\
\textstyle\sum b &= 0+0+3+0+1+1+0+0+0 = 5,\\
\textstyle\sum a^2 &= 0+0+0+9+4+4+4+4+4 = 29,\\
\textstyle\sum ab &= 0+0+0+0+2+2+0+0+0 = 4,\\
\textstyle\sum b^2 &= 0+0+9+0+1+1+0+0+0 = 11.
\end{align*}
Mean: $\vect{m}=(13/9,\ 5/9)=(1.444,\ 0.556)$. Covariance entries, in ninths and eighty-firsts:
\begin{align*}
\sigma_{aa} &= \frac{29}{9}-\left(\frac{13}{9}\right)^2 = \frac{261}{81}-\frac{169}{81} = \frac{92}{81} = 1.136,\\
\sigma_{ab} &= \frac{4}{9}-\frac{13}{9}\cdot\frac{5}{9} = \frac{36}{81}-\frac{65}{81} = -\frac{29}{81} = -0.358,\\
\sigma_{bb} &= \frac{11}{9}-\left(\frac{5}{9}\right)^2 = \frac{99}{81}-\frac{25}{81} = \frac{74}{81} = 0.914.
\end{align*}
These match \code{np.cov(..., bias=True)} on the nine pixels \tagC.
\end{worked}

\paragraph{The cost argument.} With integral images every pixel needs two boxes (outer and guard) of four
look-ups per channel. For the 31 channels of the local mean on the PRISMA frame:
\begin{align*}
\text{direct:} &\quad 3{,}280\times31\times1{,}474{,}990\approx1.50\times10^{11},\\
\text{integral:} &\quad 2\times4\times31\times1{,}474{,}990 = 365{,}797{,}520\approx3.7\times10^{8},\\
\text{ratio:} &\quad \frac{3{,}280}{2\times4} = 410.
\end{align*}
Building the tables adds about two additions per entry per channel (two running sums) over the padded grid,
$(1210+61)\times(1219+61)=1271\times1280=1{,}626{,}880$ entries:
\begin{align*}
2\times31\times1{,}626{,}880 = 100{,}866{,}560\approx1.0\times10^{8},
\end{align*}
small beside the direct method. A local covariance with $d=30$ would need $465$ product channels on top of the
30 sums and the count, 496 in all. In \code{float64} on the padded grid that is
\begin{align*}
496\times1{,}626{,}880\times8 \text{ bytes} = 6{,}455{,}459{,}840 \text{ bytes}\approx6.5\text{ GB},
\end{align*}
against $31\times1{,}626{,}880\times8=403{,}466{,}240$ bytes, about 0.4~GB, for the local mean \tagC; and every
pixel would then need its own $30\times30$ inverse. Slower, sixteen times the memory, and, as measured,
worse.

%======================================================================
\section{Eleven scores, fixed before the run}\label{sec:residualmaha:eleven}
\index{scoring methods!comparison}

On 2026-10-15 eleven scoring methods were compared on the four labelled PRISMA scenes with the F001 model at
epoch~52 (and on v0.2.0), with their parameters \emph{fixed before the run}: nothing was tuned on the
results. The comparison scripts are not in the repository; the methods are named in
\file{docs/lld/residual-mahalanobis.md}, and the numbers below are the course's record of the measured
results (IW46). A letter in brackets means ``applied on top of'': E(C) is method E run on C's score map.

\begin{drift}
F001 at epoch 52, measured 2026-10-15 (\texttt{indradhanu workings/46-the-ring-and-the-comparison.md}; the means
for C and D and the object counts 403 and 383 are confirmed by \file{docs/lld/residual-mahalanobis.md}). ROC-AUC
per scene and the mean; ``objects'' is the number of labelled objects with at least one pixel in the top 1\%.
\begin{center}\small
\begin{tabular}{@{}lcccccc@{}}
\toprule
method & set1 & set2 & set3 & set4 & mean & objects\\
\midrule
A \texttt{combined} (old default) & 0.731 & 0.755 & 0.677 & 0.698 & 0.715 & 208\\
A + SAM dark fade & 0.887 & 0.755 & 0.685 & 0.698 & 0.756 & \\
B \texttt{combined}, rank-normalised & & & & & 0.709 & \\
C residual Mahalanobis, global mean & 0.921 & 0.848 & 0.751 & 0.876 & 0.849 & \\
\textbf{D} residual Mahalanobis, ring 21/61 & \textbf{0.929} & \textbf{0.859} & \textbf{0.787} & \textbf{0.885} & \textbf{0.865} & 383\\
E(C) local z-score of the map & & & & & 0.790 & \\
F(C) difference-of-Gaussians boost & & & & & 0.718 & 403\\
G(C) $\pm1$ px shift suppression & & & & & 0.845 & \\
\bottomrule
\end{tabular}
\end{center}
Eight of the eleven are shown; the other three were chained variants, none of which beat D. Blank cells were
not recorded.
\end{drift}

\begin{figure}[htbp]\centering
\begin{tikzpicture}
\begin{axis}[xbar,width=0.82\linewidth,height=7.2cm,xmin=0.65,xmax=0.89,
  xlabel={mean ROC-AUC over the four PRISMA scenes},
  symbolic y coords={B,A,F,Af,E,G,C,D},ytick=data,
  yticklabels={{B rank-norm.},{A combined},{F(C) blob boost},{A + SAM fade},{E(C) local z},{G(C) shift supp.},{C global mean},{D ring 21/61}},
  bar width=9pt,nodes near coords,nodes near coords align={horizontal},
  every node near coord/.append style={font=\scriptsize,/pgf/number format/fixed,/pgf/number format/precision=3},
  enlarge y limits=0.08,grid=major,grid style={gray!20}]
\addplot[fill=loss!70,draw=loss!80!black] coordinates {(0.709,B) (0.715,A) (0.756,Af)};
\addplot[fill=spec!70,draw=spec!80!black] coordinates {(0.718,F) (0.790,E) (0.845,G) (0.849,C) (0.865,D)};
\end{axis}
\end{tikzpicture}
\caption{The comparison drawn to scale from the table (F001, epoch 52, 2026-10-15). Amber: scores built on the
residual's size; teal: scores built on its direction (C, D and the filters applied to C). The axis starts at
0.65 to separate the bars; the gap from A to C is $0.849-0.715=0.134$.}\label{fig:residualmaha:bars}
\end{figure}

\coursefig[0.9\linewidth]{iw-46-methods}{The same comparison as a dot plot (course figure 46.2), with object
counts where they were recorded.}{fig:residualmaha:methods}

\begin{figure}[htbp]\centering
\begin{tikzpicture}
\begin{axis}[ybar,width=0.8\linewidth,height=5.6cm,ymin=0.6,ymax=0.96,
  symbolic x coords={set1,set2,set3,set4},xtick=data,ylabel={ROC-AUC},
  bar width=8pt,legend style={at={(0.5,1.02)},anchor=south,legend columns=3,font=\small},
  grid=major,grid style={gray!20},enlarge x limits=0.18]
\addplot[fill=loss!60,draw=loss!80!black] coordinates {(set1,0.731) (set2,0.755) (set3,0.677) (set4,0.698)};
\addplot[fill=spec!40,draw=spec!80!black] coordinates {(set1,0.921) (set2,0.848) (set3,0.751) (set4,0.876)};
\addplot[fill=spec!85,draw=spec!80!black] coordinates {(set1,0.929) (set2,0.859) (set3,0.787) (set4,0.885)};
\legend{A combined, C global mean, D ring 21/61}
\end{axis}
\end{tikzpicture}
\caption{Per-scene ROC-AUC of A, C and D, to scale (F001, epoch 52). D beats C on every scene; the axis starts
at 0.6.}\label{fig:residualmaha:perscene}
\end{figure}

\paragraph{How to read it.}
\begin{itemize}
\item \textbf{Size versus direction is the big step.} Every direction-based score except the blob boost is
  above every size-based score. A to C is $0.849-0.715=0.134$.
\item \textbf{D beats C on every scene}, not only on average:
  \begin{align*}
  \text{set1: } 0.929-0.921 &= 0.008, & \text{set2: } 0.859-0.848 &= 0.011,\\
  \text{set3: } 0.787-0.751 &= 0.036, & \text{set4: } 0.885-0.876 &= 0.009.
  \end{align*}
  A gain on every scene is more convincing than a gain in the mean.
\item \textbf{Spatial post-filters did not help AUC.} A local z-score (each score re-expressed against its
  neighbourhood's spread), a blob boost (a difference of two blurs that favours compact blobs) and shift
  suppression (meant to discount errors that vanish if the reconstruction moves by one pixel) all scored below
  C.
\item \textbf{AUC and object counts can disagree.} F(C) catches the most objects, 403 against D's 383, but its
  AUC is 0.718. AUC judges the ranking of \emph{every} pixel; the object count looks only at the top 1\% and
  gives one point per object touched (\cref{ch:measures}). A blob boost lifts compact objects into the top 1\%
  while scrambling the ranking everywhere else. Against D it gave up $0.865-0.718=0.147$ of AUC for
  $403-383=20$ more objects. The product needs the ranking: an analyst moves the threshold, so the whole order
  matters.
\end{itemize}
D became Indradhanu's default on 2026-10-15 (roadmap row OB-13): the resolver lists
\code{residual_mahalanobis} first for both Indradhanu models, and \code{combined}, \code{L1} and \code{SAM}
stay selectable through the per-model override (\file{backend/allotrope/foundation_models/resolver.py}, lines
152--173).

\begin{drift}
The same comparison with the production model v0.2.0 (\file{research/experiments/dual-window-mahalanobis/README.md},
section 3): \texttt{combined} 0.662, L1 0.702, SAM 0.570, MNF-RX 0.752, PCA30-RX on the cube 0.831, residual
Mahalanobis global mean 0.839, ring 21/61 0.854. Set~3 is the exception there: PCA30-RX on the cube (0.766)
beats residual Mahalanobis with the global mean (0.709).
\end{drift}

\subsection{Three honest caveats}

\paragraph{The window was chosen on these scenes.} The LLD says it plainly: ``The window was the best of four
in the original sweep on the same four scenes, so the margin is slightly flattering.'' When you pick the best
of four settings on a test, part of its lead is luck that will not repeat on new scenes. D's lead over C,
$0.865-0.849=0.016$, is small enough that some of it may be that luck. That is why \code{local_window=None},
the global version, which has no window to choose, stays one argument away.

\paragraph{Four scenes are very few.}\index{sign test} If the two methods were truly equal, each scene would be a coin flip for
which one wins:
\begin{align*}
P(\text{all 4 go to D}) &= \left(\tfrac12\right)^4 = \tfrac{1}{16},\\
P(\text{all 4 go one way or the other}) &= 2\times\tfrac{1}{16} = \tfrac{2}{16} = 0.125.
\end{align*}
Even a clean sweep happens one time in eight by chance, so no paired test on four scenes can reach the usual
$0.05$. The evidence is the consistency and the size of the gap, not a significance test (\cref{ch:bootstrap}
develops a better tool).

\paragraph{The ground truth may sit about one pixel off.} AUC counts that against every method equally; the
object count allows one pixel of slack.

\begin{sourcenote}[: the class-aware variant]
The same module also defines \code{class_aware_residual_mahalanobis} (\file{residual_mahalanobis.py}, lines
75--102): a separate whitening per surface class from a segmentation, with same-class rings, and classes under
5,000 pixels merged into land. The product's scoring action does not call it; it is used by the evaluator's
scorers (\file{app/evaluation/scorers.py}, lines 32--59) for experiments. When this book says ``residual
Mahalanobis'' it means the default, class-blind version.
\end{sourcenote}

%======================================================================
\begin{confusions}
\confusion{Is residual Mahalanobis a new model?}{No. The model and its reconstruction are unchanged; only the
step that turns a residual into a score differs. A checkpoint needs no retraining to use it.}
\confusion{Plain RX was dropped for exploding. Why is this allowed back?}{Two changes: only the 30 largest
directions are used, so the tiny, noisy eigenvalues that exploded are discarded; and it runs on the residual,
which is closer to one cloud of ordinary errors than the raw cube.}
\confusion{Is a high score a big error?}{Not necessarily, and that is the point. $\vect{Q}$ in
\cref{we:residualmaha:three} has the smallest error and the highest score.}
\confusion{There is already a \texttt{mahalanobis} in the product.}{That is the classical detector on the cube
(MNF-RX, thermal GRX). The new method is deliberately named \texttt{residual\_mahalanobis} so the two are never
confused in configurations or the interface.}
\confusion{Local mean, so the covariance is local too?}{No. Only the mean comes from the ring. The directions
and spreads are the whole scene's.}
\confusion{Why does a ring pixel's validity matter if its whitened value is already 0?}{A zero value would
still be counted in the denominator. The count channel makes the mean divide by valid pixels only.}
\confusion{Can I threshold the score at, say, 10?}{No. Scores are squared distances with a heavy right tail, and
their scale depends on the scene. Like every score in the product, it is thresholded by percentile
(\cref{ch:material}) and drawn with a square-root stretch.}
\end{confusions}

\begin{checkyourself}
\item\label{ex:residualmaha:1} Using the four ordinary pixels of \cref{we:residualmaha:three}, score a new pixel
  with residual $(1,1,2)$ by $\Cov_r\inv$, and again by rotate-and-divide. Compare its L1 with that of
  $\vect{Q}$.
\item\label{ex:residualmaha:2} Using the integral image of \cref{we:residualmaha:integral}, compute the sum of
  value rows 0--1, columns 2--3, and the toy ring mean (3$\times$3 outer, 1$\times$1 guard) around value row 1,
  column 2. Check both by adding values.
\item\label{ex:residualmaha:3} Four pixels in two channels have $\sum a=8$, $\sum b=4$, $\sum a^2=20$,
  $\sum ab=6$, $\sum b^2=6$. Compute the mean and the covariance (divide by $n$) from these sums alone.
\item\label{ex:residualmaha:4} For a window \code{(11, 31)}: how many pixels are in the ring, what is the
  fallback threshold if $d=20$, and by what factor do integral images reduce the per-pixel cost of the ring
  sums?
\item\label{ex:residualmaha:5} In the 3-band example keep only the top two directions. Score $(1,1,2)$ and
  $\vect{Q}$, and say in one sentence what truncation gave up.
\item\label{ex:residualmaha:6} A fifth labelled scene arrives and D beats C there too. What is now the chance
  of a clean sweep in one direction or the other if the methods were equal?
\end{checkyourself}

\nextchapter{ch:material}
